\documentclass{includes/Thesis}

\addbibresource{parts/refs.bib}

\begin{thesis}[
    institute = ПРАВИТЕЛЬСТВО РОССИЙСКОЙ ФЕДЕРАЦИИ\\ФЕДЕРАЛЬНОЕ ГОСУДАРСТВЕННОЕ АВТОНОМНОЕ\\ОБРАЗОВАТЕЛЬНОЕ УЧРЕЖДЕНИЕ ВЫСШЕГО ОБРАЗОВАНИЯ\\НАЦИОНАЛЬНЫЙ ИССЛЕДОВАТЕЛЬСКИЙ УНИВЕРСИТЕТ\\<<ВЫСШАЯ ШКОЛА ЭКОНОМИКИ>>,
    faculty = Санкт-Петербургская школа физико-математических\\и компьютерных наук НИУ ВШЭ --- Санкт-Петербург\\Образовательная программа <<Прикладная математика и информатика>>,
    year = 2026,
    city = Санкт-Петербург,
    %
    title = Улучшенная нижняя оценка для константы Чватала--Санкоффа методом динамического программирования с окнами,
    %
    authorGroup = БПМ222С,
    authorName = Д.~Е.~Юкачев,
    %
    academicTeacherTitle = Старший преподаватель,
    academicTeacherName = С.~В.~Копелиович,
    %
    academicSupervisorTitle = {Академический руководитель\\образовательной программы\\<<Прикладная математика и информатика>>,\\Доцент, канд. техн. наук},
    academicSupervisorName = А.~И.~Храбров,
    %
    isAcademic,
    UDC = 004.421,
    educationalProgram = 01.03.02 <<Прикладная математика и информатика>>,
    %
    keywordsRu = константа Чватала--Санкоффа; наибольшая общая подпоследовательность; динамическое программирование; нижние оценки; случайные строки; задача Улама; Monte-Carlo,
    keywordsEn = Chv\'atal--Sankoff constant; longest common subsequence; dynamic programming; lower bounds; random strings; Ulam problem; Monte-Carlo,
]

    \setAbstractResource{parts/abstract-ru}{parts/abstract-en}

    \setTerminologyResource{parts/terminology}

    \setIntroResource{parts/intro}

    \addChapter{Предметная область и обзор литературы}{parts/chapter1}
    \addChapter{Метод оконного динамического программирования}{parts/chapter2}
    \addChapter{Корректность и вычислительная сложность}{parts/chapter3}
    \addChapter{Реализация и история оптимизаций}{parts/chapter4}
    \addChapter{Численные эксперименты и верификация}{parts/chapter5}

    \setConclusionResource{parts/conclusion}

    \addAppendix{Список сокращений и обозначений}{parts/appendix-a}
    \addAppendix{Параметры экспериментов}{parts/appendix-b}
    \addAppendix{Ключевые фрагменты исходного кода}{parts/appendix-c}
    \addAppendix{Дополнительные таблицы}{parts/appendix-d}

\end{thesis}
